\chapter{Ohne Ende}

\section{Der Fall}

\begin{merksatz}
\boxtitel{Wie oft soll man untersuchen?}
Eine Vorsorgeuntersuchung entdeckt eine Erkrankung, bevor sie fortschreitet. Wie groß soll der Abstand sein?

\begin{center}
{\small
\renewcommand{\arraystretch}{1.5}
\begin{tabular}{@{}P{3.0cm}P{2.4cm}@{}}
\toprule
Neuerkrankung je Jahr & $2\,\%$\\
Fortschreiten je Jahr & $25\,\%$\\
Entdeckungsrate & $90\,\%$\\
Kosten je Untersuchung & $100$ €\\
\bottomrule
\end{tabular}}
\end{center}

\schritt{Simuliert über dreißig Jahre}
\begin{center}
{\small
\renewcommand{\arraystretch}{1.48}
\begin{tabular}{@{}P{2.4cm}P{2.2cm}P{2.2cm}P{2.2cm}@{}}
\toprule
\textbf{Abstand} & \textbf{QALY} & \textbf{Kosten} & \textbf{je QALY}\\
\midrule
$10$ Jahre & $15{,}735$ & $199$ € & ---\\
$5$ Jahre & $16{,}084$ & $361$ € & $464$ €\\
$3$ Jahre & $16{,}365$ & $584$ € & $794$ €\\
$2$ Jahre & $16{,}587$ & $868$ € & $1279$ €\\
$1$ Jahr & $16{,}833$ & $1722$ € & $3472$ €\\
\bottomrule
\end{tabular}}
\end{center}
\kursiv{Die letzte Spalte ist der Zuwachs gegenüber der Zeile darüber.}

\begin{center}
\begin{minipage}{0.9\textwidth}
\textbf{Jede Verkürzung bringt etwas --- und jede kostet mehr als die vorige.}
\end{minipage}
\end{center}

\begin{center}
\begin{minipage}{0.88\textwidth}
\kursiv{\textbf{Es gibt kein \enquote{richtiges} Intervall aus der Rechnung allein.} Es gibt eine Folge von Stufen mit steigenden Preisen, und welche man noch nimmt, hängt davon ab, was ein gewonnenes Jahr wert sein soll. Kapitel~10 macht daraus eine Grenze, Kapitel~14 setzt einen Preis an.}
\end{minipage}
\end{center}
\end{merksatz}

\section{Die Gleichung ohne Endzeit}

\begin{merksatz}
\boxtitel{Satz 8.1 --- der Fixpunkt}
Bei unbegrenztem Zeithorizont und Diskontfaktor $\gamma < 1$ hat
\[ V(z) = \max_a \left\{ u(z,a) + \gamma \sum_{z'} \Pr[z'\mid z,a]\,V(z') \right\} \]
genau eine Lösung, und die Wertiteration
\[ V^{(k+1)} = \mathcal{B}\,V^{(k)} \]
konvergiert von jedem Startwert aus gegen sie.

\schritt{Beweis}
\kursiv{Erstens --- die Abbildung zieht zusammen:} Für zwei Wertfunktionen $V$, $W$ gilt an jeder Stelle
\[ \bigl|\mathcal{B}V(z) - \mathcal{B}W(z)\bigr| \ \le\ \gamma\,\max_{z'}\bigl|V(z')-W(z')\bigr| , \]
denn das Maximum zweier Ausdrücke unterscheidet sich höchstens so stark wie die Ausdrücke selbst, und die Übergangswahrscheinlichkeiten summieren sich zu eins.

\kursiv{Zweitens --- Banach:} Eine Kontraktion mit Faktor $\gamma < 1$ auf einem vollständigen Raum hat genau einen Fixpunkt, und die Iteration konvergiert gegen ihn --- Band 3, Kapitel~9.

\kursiv{Drittens --- die Geschwindigkeit:} Der Abstand schrumpft je Schritt um den Faktor $\gamma$. $\blacksquare$

\begin{center}
{\small
\renewcommand{\arraystretch}{1.46}
\begin{tabular}{@{}P{2.4cm}P{2.4cm}P{2.6cm}@{}}
\toprule
$\gamma$ & \textbf{Schritte für $1\,\%$} & \textbf{entspricht}\\
\midrule
$0{,}90$ & $44$ & $10$ Jahre Sicht\\
$0{,}97$ & $151$ & $33$ Jahre\\
$0{,}99$ & $458$ & $100$ Jahre\\
\bottomrule
\end{tabular}}
\end{center}

\begin{center}
\begin{minipage}{0.88\textwidth}
\kursiv{\textbf{Der Diskontfaktor ist zugleich die Rechenbremse und der Zeithorizont}: $\gamma = 0{,}97$ entspricht der Diskontrate von drei Prozent aus Kapitel~6 und einer wirksamen Sicht von etwa dreiunddreißig Jahren. Wer weiter blicken will, muss länger rechnen --- und braucht Daten, die so weit reichen.}
\end{minipage}
\end{center}
\end{merksatz}

\section{Der Bruchfall}

\begin{warnung}
\boxtitelr{Was die schöne Theorie verlangt}
\begin{center}
{\small
\renewcommand{\arraystretch}{1.52}
\begin{tabular}{@{}P{3.2cm}P{5.4cm}@{}}
\toprule
\textbf{Bekannte Über\-gänge} & Satz 8.1 setzt sie als gegeben voraus. In der Medizin sind sie geschätzt --- aus Daten, die selbst unter der bisherigen Strategie entstanden.\\
\textbf{Zustandsraum} & Jede zusätzliche Größe vervielfacht ihn. Der Fluch der Dimension trifft hier härter als beim Lernen.\\
\textbf{Verlässliche Nutzen} & Der Wert wird über Jahrzehnte aufsummiert; kleine Fehler in $u$ wirken groß.\\
\bottomrule
\end{tabular}}
\end{center}

\begin{center}
\begin{minipage}{0.88\textwidth}
\kursiv{Die erste Zeile enthält die Falle, die Band 23 in Kapitel~12 beschrieb: \textbf{Die beobachteten Übergänge gelten für die Strategie, unter der beobachtet wurde} --- nicht für die vorgeschlagene. Eine berechnete Strategie, die deutlich vom Bisherigen abweicht, bewegt sich damit in einem Bereich ohne Daten. Kapitel~11 behandelt Verfahren, die darauf Rücksicht nehmen.}
\end{minipage}
\end{center}
\end{warnung}

\section{Übungsaufgaben}

\begin{aufgabe}
\boxtitel{Aufgabe 8.1 --- Eine Wertiteration von Hand}
Zwei Zustände: gesund ($u = 1$) und krank ($u = 0{,}5$). Aus gesund wird man mit $10$ Prozent krank; aus krank wird man mit Behandlung ($u$-Abzug $0{,}1$) zu $40$ Prozent gesund, ohne zu $10$. $\gamma = 0{,}9$.

\loesung
\schritt{Schritt 1 --- Start}
\[ V^{(0)}(\text{g}) = V^{(0)}(\text{k}) = 0 \]

\schritt{Schritt 2 --- Erste Iteration}
\[ V^{(1)}(\text{g}) = 1 + 0{,}9\cdot0 = 1{,}000 \]
\[ V^{(1)}(\text{k}) = \max\{0{,}5;\ 0{,}4\} = 0{,}500 \quad (\text{ohne Behandlung}) \]

\schritt{Schritt 3 --- Zweite Iteration}
\[ V^{(2)}(\text{g}) = 1 + 0{,}9\,[0{,}9\cdot1{,}000 + 0{,}1\cdot0{,}500] = 1 + 0{,}9\cdot0{,}950 = 1{,}855 \]
\[ V^{(2)}(\text{k}, \text{behandeln}) = 0{,}4 + 0{,}9\,[0{,}4\cdot1{,}000+0{,}6\cdot0{,}500] = 0{,}4+0{,}9\cdot0{,}700 = 1{,}030 \]
\[ V^{(2)}(\text{k}, \text{nicht}) = 0{,}5 + 0{,}9\,[0{,}1\cdot1{,}000+0{,}9\cdot0{,}500] = 0{,}5+0{,}9\cdot0{,}550 = 0{,}995 \]
\[ V^{(2)}(\text{k}) = \mathbf{1{,}030} \quad (\text{behandeln}) \]

\begin{center}
\begin{minipage}{0.9\textwidth}
\textbf{Die Entscheidung kippt zwischen der ersten und der zweiten Iteration.}
\end{minipage}
\end{center}

\schritt{Schritt 4 --- Warum}
\begin{center}
\begin{minipage}{0.88\textwidth}
\kursiv{\textbf{In der ersten Iteration zählt nur das Jetzt, und da ist Nichtstun angenehmer.} Erst wenn die Zukunft mitgerechnet wird, lohnt der sofortige Abzug von $0{,}1$, weil er die Aussicht auf den besseren Zustand vervierfacht. Kurzsichtigkeit ist also kein Fehler des Verfahrens, sondern der Grenzfall $\gamma = 0$ --- und beschreibt gut, was unter Zeitdruck geschieht.}
\end{minipage}
\end{center}

\schritt{Schritt 5 --- Der Grenzwert}
Nach weiteren Iterationen: $V(\text{g}) = 9{,}02$, $V(\text{k}) = 7{,}93$.

\kursiv{Probe an beiden Gleichungen:}
\[ 1 + 0{,}9\,[0{,}9\cdot9{,}02 + 0{,}1\cdot7{,}93] = 1 + 0{,}9\cdot8{,}911 = 9{,}02 \quad \checkmark \]
\[ 0{,}4 + 0{,}9\,[0{,}4\cdot9{,}02 + 0{,}6\cdot7{,}93] = 0{,}4 + 0{,}9\cdot8{,}366 = 7{,}93 \quad \checkmark \]
\end{aufgabe}

\begin{aufgabe}
\boxtitel{Aufgabe 8.2 --- Das Intervall wählen}
Ab welchem Preis je gewonnenem Lebensjahr lohnt sich welches Untersuchungsintervall?

\loesung
\schritt{Schritt 1 --- Die Stufen aus Abschnitt 8.1}
\begin{center}
{\small
\renewcommand{\arraystretch}{1.48}
\begin{tabular}{@{}P{2.4cm}P{2.6cm}P{2.6cm}@{}}
\toprule
\textbf{Übergang} & \textbf{Preis je QALY} & \textbf{lohnt ab}\\
\midrule
$10 \to 5$ & $464$ € & $464$ €\\
$5 \to 3$ & $794$ € & $794$ €\\
$3 \to 2$ & $1279$ € & $1279$ €\\
$2 \to 1$ & $3472$ € & $3472$ €\\
\bottomrule
\end{tabular}}
\end{center}

\schritt{Schritt 2 --- Die Regel}
\begin{center}
\begin{minipage}{0.9\textwidth}
\textbf{Man geht so lange eine Stufe weiter, wie der Preis der Stufe unter der Zahlungsbereitschaft liegt.}
\end{minipage}
\end{center}

\schritt{Schritt 3 --- Bei $40\,000$ Euro je QALY}
Alle vier Stufen lohnen --- das jährliche Intervall gewinnt.

\schritt{Schritt 4 --- Der Haken}
\begin{center}
\begin{minipage}{0.88\textwidth}
\kursiv{\textbf{Das Modell kennt keine Schäden des Untersuchens} --- keine falsch positiven Befunde, keine Folgediagnostik, keine Überdiagnose. Nimmt man sie auf, so kehren sich die späten Stufen um: Der zusätzliche Nutzen sinkt, der zusätzliche Schaden steigt, und ab einem gewissen Punkt wird häufigeres Untersuchen \kursiv{schädlich}, nicht nur teuer. Genau darüber wird bei Krebsvorsorgeprogrammen gestritten --- und der Streit ist einer über Modellstruktur, nicht über Zahlen.}
\end{minipage}
\end{center}
\end{aufgabe}

\begin{aufgabe}
\boxtitel{Aufgabe 8.3 --- Wie weit reicht der Blick?}
Welchem Zeithorizont entspricht ein Diskontfaktor $\gamma$?

\loesung
\schritt{Schritt 1 --- Die Summe}
\[ \sum_{t=0}^{\infty} \gamma^{t} = \frac{1}{1-\gamma} \]

\schritt{Schritt 2 --- Die Werte}
\begin{center}
{\small
\renewcommand{\arraystretch}{1.46}
\begin{tabular}{@{}P{2.0cm}P{2.4cm}P{2.6cm}@{}}
\toprule
$\gamma$ & $1/(1-\gamma)$ & \textbf{Diskontrate}\\
\midrule
$0{,}90$ & $10$ & $11{,}1\,\%$\\
$0{,}95$ & $20$ & $5{,}3\,\%$\\
$0{,}97$ & $33$ & $3{,}1\,\%$\\
$0{,}99$ & $100$ & $1{,}0\,\%$\\
\bottomrule
\end{tabular}}
\end{center}

\kursiv{Probe:} $\gamma = 1/(1+d)$, also $d = 1/\gamma - 1 = 1/0{,}97-1 = 0{,}0309$.

\schritt{Schritt 3 --- Die Deutung}
\begin{center}
\begin{minipage}{0.88\textwidth}
\kursiv{\textbf{Der Kehrwert von $1-\gamma$ ist die Zahl der Jahre, die praktisch zählen.} Eine Entscheidung mit $\gamma = 0{,}90$ ignoriert alles jenseits eines Jahrzehnts --- was bei einer $78$-jährigen Patientin vertretbar sein mag und bei einem Kind unvertretbar ist. In technischen Anwendungen des bestärkenden Lernens wird $\gamma$ oft nach der Rechengeschwindigkeit gewählt; in der Medizin ist es eine Wertentscheidung.}
\end{minipage}
\end{center}
\end{aufgabe}

\begin{geschichte}
\boxtitelg{Von der Lagerhaltung zur Leitlinie}
Markow-Entscheidungsprozesse entstanden in den 1950er Jahren --- \Person{Ronald Howard} schrieb 1960 das erste Lehrbuch.

\begin{center}
\begin{minipage}{0.9\textwidth}
\kursiv{Die frühen Anwendungen betrafen Lagerhaltung, Instandhaltung und Warteschlangen. \textbf{\Person{Lloyd Shapley} hatte das Verfahren bereits 1953 für stochastische Spiele formuliert} --- dieselbe Person, die in Band 21 den Shapley-Wert beisteuerte, und derselbe Gedanke, nur mit einem Gegenspieler statt der Natur.}
\end{minipage}
\end{center}

\schritt{Und die Wiederkehr}
\begin{center}
\begin{minipage}{0.9\textwidth}
\kursiv{Das bestärkende Lernen der 1990er Jahre ist derselbe Rahmen ohne Kenntnis der Übergänge: \textbf{Man lernt sie aus Erfahrung, statt sie vorauszusetzen.} Der Erfolg bei Brett- und Videospielen hat das Verfahren populär gemacht; die Übertragung auf die Medizin steht vor dem Problem, dass man dort nicht beliebig oft probieren kann.}
\end{minipage}
\end{center}
\end{geschichte}

\begin{forschung}
\boxtitelp{Strategien heute}
\begin{center}
{\small
\renewcommand{\arraystretch}{1.45}
\begin{tabular}{@{}P{3.0cm}P{5.2cm}@{}}
\toprule
\textbf{Teilweise sichtbar} & Der Zustand ist unbekannt; der Glaube über ihn wird zum Zustand.\\
\textbf{Sicheres Lernen} & Strategien, die den beobachteten Bereich nicht verlassen.\\
\textbf{Erklärbare Regeln} & Einfache Schwellenregeln statt optimaler, aber unverständlicher Strategien.\\
\bottomrule
\end{tabular}}
\end{center}

\schritt{Zur dritten Zeile}
\begin{center}
\begin{minipage}{0.88\textwidth}
\kursiv{Die berechnete beste Strategie ist oft eine Tabelle ohne erkennbares Muster. \textbf{Eine um wenige Prozent schlechtere Strategie, die sich als \enquote{behandle ab Wert $x$} aussprechen lässt, wird eher befolgt} --- und eine befolgte schlechtere Strategie schlägt eine ignorierte bessere. Das ist dieselbe Abwägung wie zwischen Baum und Wald in Band 22, Kapitel~9.}
\end{minipage}
\end{center}
\end{forschung}

\begin{merksatz}
\boxtitel{Was Sie jetzt können}
\begin{center}
\renewcommand{\arraystretch}{1.5}
\begin{tabular}{@{}P{3.0cm}P{5.6cm}@{}}
\toprule
\textbf{Satz 8.1} & Kontraktion, ein Fixpunkt\\
\textbf{Wertiteration} & konvergiert mit Faktor $\gamma$\\
\textbf{$1/(1-\gamma)$} & die Zahl der zählenden Jahre\\
\textbf{Stufen} & jede Verkürzung teurer als die vorige\\
\textbf{Der Bruchfall} & Übergänge gelten für die alte Strategie\\
\bottomrule
\end{tabular}
\end{center}

\begin{center}
\begin{minipage}{0.88\textwidth}
\kursiv{Bisher waren die Wahrscheinlichkeiten bekannt. Das nächste Kapitel handelt vom Fall, in dem man sie erst durch Handeln erfährt --- und dabei Patienten einsetzt.}
\end{minipage}
\end{center}
\end{merksatz}
